\section{Economic accuracy of approximately optimized continuations}\label{sec:actionwise82}
A qualification tolerance need not preserve every acquired action to preserve an economic conclusion. The bound in Theorem~\ref{thm:stability81} charges the full cost range on each uncertified cell. We now retain the particular actions that remain possible and the direction of their economic effects. This gives a cost guarantee for an entire optimization sublevel, including approximately optimal readouts, without treating an inconclusive switch certificate as an observed policy failure.

Fix a finite first-date acquisition partition $\mathcal C$, its initial-law masses $p_C$, a robust finite action menu on each cell, and a common feasible continuation policy $\tau$ after that date. These objects are held fixed throughout the comparison. The horizon may be any finite horizon for which the conditional costs below are enclosed. Use the fixed-label qualification objective $F$, attained readout $\widehat\theta$, value $U=F(\widehat\theta)$, reference action $a_C$, and switch lower bounds $D_{Cb}$ from Section~\ref{sec:calibration81}. For $\epsilon\ge0$, define
\begin{equation}\label{eq:possible-actions82}
 \mathcal A_C(\epsilon)=\{a_C\}\cup
       \{b\ne a_C:D_{Cb}\le U+\epsilon\}.
\end{equation}
The weak inequality deliberately retains ties. Let $\pi_\theta$ use a minimizing critic action at each cell representative and then $\tau$. Write $Q^\tau(x,b)=c_0(x,b)+\beta_0P_0^bJ_1^\tau(x)$, and, on every positive-mass cell, define the original-law contrast
\[
 \Delta_C(b)=\mathbb E[Q^\tau(X_0,b)-Q^\tau(X_0,a_C)\mid X_0\in C].
\]
Suppose a finite calculation provides $\ell_{Cb}\le\Delta_C(b)\le u_{Cb}$ for every retained action, with both endpoints zero for $a_C$. Set
\begin{equation}\label{eq:class-band82}
 \underline v_\epsilon=\sum_Cp_C\min_{b\in\mathcal A_C(\epsilon)}\ell_{Cb},
 \qquad
 \overline v_\epsilon=\sum_Cp_C\max_{b\in\mathcal A_C(\epsilon)}u_{Cb}.
\end{equation}
Zero-mass cells make no contribution.

\begin{theorem}[Actionwise optimization-to-cost transfer]\label{thm:actionwise82}
Under the preceding conditions, every feasible $\theta$ satisfying $F(\theta)\le U+\epsilon$ obeys
\begin{equation}\label{eq:uniform-cost82}
 \underline v_\epsilon\le J^{\pi_\theta}-J^{\pi_{\widehat\theta}}
                         \le\overline v_\epsilon.
\end{equation}
In particular the conclusion holds for every $\epsilon$-optimal readout. If all complete costs are in $[0,H]$ and the actionwise intervals are intersected with $[-H,H]$, then
\[
 -Hq_\epsilon\le\underline v_\epsilon\le0\le
       \overline v_\epsilon\le Hq_\epsilon,
\]
where $q_\epsilon$ is the initial mass of cells retaining a nonreference action.
For two such optimizer classes, an attained-policy simultaneous interval $[l_{jk},u_{jk}]$ gives, on the same event, the classwide interval
\begin{equation}\label{eq:class-pair82}
 [\,l_{jk}+\underline v_j-\overline v_k,\quad
       u_{jk}+\overline v_j-\underline v_k\,].
\end{equation}
These conclusions require no probability bound over the readouts in the sublevels.
\end{theorem}

\begin{proof}
Any critic-minimizing action $b\ne a_C$ satisfies $g_{Cb}(\theta)\le0$, so $D_{Cb}\le F(\theta)\le U+\epsilon$. It therefore belongs to \eqref{eq:possible-actions82}. Conditioning on the initial cell expresses the policy-cost difference as $\sum_Cp_C\Delta_C(b_C(\theta))$, which is between the endpoints in \eqref{eq:class-band82}. The infimum of $F$ is at most $U$, giving the optimization conclusion. A cell retaining only $a_C$ contributes zero; the signed support bounds every other contribution. Finally add and subtract the two attained costs to obtain \eqref{eq:class-pair82}. All steps after the attained-policy event are deterministic.
\end{proof}

The cellwise extrema relax the common-coefficient restriction and need not be attained by one readout. The theorem is consequently not a sharp joint characterization of the optimizer class. Its useful feature is different: it can retain an identified economic comparison when some actions remain uncertified and the binary $Hq_\epsilon$ allowance is too large. The original Bellman optimum and the constructive accuracy results are unchanged. A reference expected-loss upper bound can also be enlarged by $\overline v_\epsilon$; the actionwise comparison itself is not a new bound obtained by replacing the original optimum with a menu optimum.

For the executed two-date economy, the conditional contrasts are computed from the original stage costs, transitions and fixed terminal actor. Common innovations preserve their actual marginal laws. Centered polynomial differences retain cancellation, and rectangle queries enclose every crossed acquired cell. The final innovation is integrated analytically. Appendix~\ref{app:actionwise82} proves whole-bin validity, deterministic refinement convergence under vanishing arithmetic error, and the work bound. It also separates integration accuracy from the independent statistical interval to which \eqref{eq:class-pair82} is applied.
